A scene graph generation method reports a ten-point gain in mean recall; a long-tailed
classifier reports two points of accuracy. What did the gain buy? On benchmarks whose labels
are partly determined by a feature both compared models observe --- the object-class pair in
relation prediction, the class frequency in long-tailed recognition --- such a gain mixes
two improvements: the new model may fit the label frequencies \emph{within} groups of that
feature more closely, which a frequency table also achieves and a shift in those frequencies
removes, or it may match its predictions to the individual cases they were made for. Mean
recall, zero-shot recall and significance tests all leave the two mixed. We decompose the
gain exactly. Score each model's prediction for one test case against the label of a
\emph{different} case in the same group: averaging over such relabellings gives the
\emph{transported gain}, and what remains is a \emph{within-group covariance gain}. The
identity is exact in finite samples, needs no fitted baseline, no held-out fold and no tuning
parameter, and costs $O(NK)$, so any pair of frozen models can be audited from cached
predictions. The remainder is an order-two U-statistic, holds for every Bregman score, and
admits cluster-robust interval estimation. Applied to two released MOTIFS checkpoints on
VG150 with the original authors' code, a ten-point mean-recall improvement corresponds to a
proper-score difference that is overwhelmingly transported, with a covariance remainder an
order of magnitude smaller under either of two proper scores. The unmodified estimator
transfers to a frozen-CLIP variant and to long-tailed image and text classification,
exposing compositions that aggregate scores hide --- including a model that loses on every
aggregate measure while discriminating individual cases significantly better than the model
it appears to lose to.
